\chapter{SOLUTIONS OF LINEAR DIFFERENTIAL EQUATIONS}

The irrationality of $\frac{\tan a}{a}$ for rational $a^2 \neq 0$, which includes the irrationality of $\pi$, was discovered by Lambert nearly two hundred years ago. Lambert's work was generalized by Legendre who considered the power series
\[
y = f_{\alpha}(x) = \sum_{n=0}^{\infty} \frac{x^n}{n! \alpha(\alpha+1) \cdots (\alpha+n-1)} \quad (\alpha \neq 0, -1, -2, \ldots)
\]
satisfying the linear differential equation of second order
\[xy'' + \alpha y' = y.\]
He obtained the continued fraction expansion
\[\frac{y}{y'} = \alpha + \frac{x}{\alpha + 1 + \frac{x}{\alpha + 2 + \cdots}}\]
and proved the irrationality of $y/y'$ for all rational $x \neq 0$ and all rational $\alpha \neq 0, -1, -2, \ldots.$ In the special case $\alpha = \frac{1}{2}$ we have
\[y = \cosh (2\sqrt{x}), \quad y' = \sinh (2\sqrt{x})/\sqrt{x},\]
so that Legendre's theorem contains the irrationality of $\frac{\tan a}{a}$ for rational $a^2 \neq 0$. In more recent times, Stridsberg proved the irrationality of $y$ and of $y'$, separately, for rational $x \neq 0$ and rational $\alpha \neq 0, -1, \ldots$, and Maier showed that neither $y$ nor $y'$ is a quadratic irrationality. Maier's work suggested the idea of introducing more general approximation forms which enabled me to prove that the numbers $y$ and $y'$ are not connected by any algebraic equation with algebraic coefficients, for any algebraic $x \neq 0$ and any rational $\alpha \neq 0, \pm \frac{1}{2}, -1, \pm \frac{3}{2}, \ldots.$ The excluded case of an integer $+ \frac{1}{2}$ really is an exception, since then the function $f_{\alpha}(x)$ satisfies an algebraic differential equation of first order whose coefficients are polynomials in $x$ with rational numerical coefficients; this follows from the explicit formulas
\[f_{k+\frac{1}{2}} = \frac{1}{2} \cdot \frac{3}{2} \cdot \cdots \cdot \left(k - \frac{1}{2}\right) D^{k} \cosh (2\sqrt{x}),\]
\[f_{\frac{1}{2}-k} = \frac{(-1)^{k} x^{k+\frac{1}{2}}}{\frac{1}{2} \cdot \frac{3}{2} \cdot \cdots \cdot \left(k - \frac{1}{2}\right)} D^{k+1} \sinh (2\sqrt{x})\]
\[(k=0,1,2,\ldots).\]
For instance, in case $\alpha = \frac{1}{2}$, the differential equation is
\[y^2 - xy'^2 = 1.\]
In the excluded case, however, Lindemann's theorem shows that $y$ and $y'$ are both transcendental for any algebraic $x \neq 0$.

We are now going to develop a general method for transcendency proofs involving solutions of linear differential equations and to apply it to the function $y=f_{\alpha}(x)$.

\section{Functions of type E}

A function $y = f(x)$ is called of type E, or an E-function, if
\[f(x) = \sum_{n=0}^{\infty} c_n \frac{x^n}{n!}\]
is a power series satisfying the following three conditions:

1) All coefficients $c_n$ belong to the same algebraic number field of finite degree over the rational number field.

2) If $\epsilon$ is any positive number, then $\overline{|c_n|} = O(n^{n\epsilon})$, as $n \to \infty$.

3) There exists a sequence $q_0, q_1, \ldots$ of positive rational integers such that $q_n c_k$ is integral for $k = 0, 1, \ldots, n$ and $n = 0, 1, 2, \ldots$, and that $q_n = O(n^{n\epsilon})$.

We shall recall to memory that the symbol $\overline{|c_n|}$ designates the maximum of the absolute values of all conjugates of the algebraic number $c_n$. The second condition means that $\overline{|c_n|}$, as a function of $n$, increases less rapidly than any positive power of $n^n$. This implies in particular, that $y$ is an entire function of $x$. The third condition states that the least common rational denominator of $c_0, \ldots, c_n$ increases less rapidly than any positive power of $n^n$.

Obviously, any polynomial with algebraic coefficients and the exponential function $e^x$ are examples of E-functions.

It is clear that the derivative $y'$ of an E-function again is an E-function. It is also clear that the sum of two E-functions is an E-function. The same is true for the product: If also
\[g(x) = \sum_{n=0}^{\infty} d_n \frac{x^n}{n!}\]
is an E-function and $r_n$ denotes the least positive rational integral common denominator of $d_0, d_1, \ldots, d_n$, then
\[
\begin{split} f(x)g(x) &= \sum_{n=0}^{\infty} e_n \frac{x^n}{n!}, \\ e_n &= \sum_{k=0}^{n} \binom{n}{k} c_k d_{n-k}, \end{split}
\]
so that
\[\overline{|e_n|} \le (1+1)^n \max_{k \le n} \overline{|c_k d_{n-k}|} = 2^n O(n^{2n\epsilon}) = O(n^{3n\epsilon}),\]
and the positive rational integer
\[q_n r_n = O(n^{2n\epsilon})\]
is a common denominator of $e_0, \ldots, e_n$. This shows that the E-functions constitute a ring. Finally, also $f(ax)$ is of type E, for any algebraic constant $a$.

Later we shall study E-functions $y_1 = E_1, \ldots, y_m = E_m$ which satisfy a system of homogeneous linear differential equations of the first order
\begin{equation}
y'_{k} = \sum_{l=1}^{m} Q_{kl}(x) y_{l} \quad (k=1,\ldots,m) \tag{49}\label{eq:49}
\end{equation}
whose coefficients $Q_{kl}$ are rational functions of $x$. Let the polynomial $T(x)$ be the least common denominator of the $m^2$ rational functions $Q_{kl}(x)$ and consider the numerical coefficients of $T(x)$ and the $T(x) Q_{kl}(x)$ as unknown quantities. If we insert the power series $E_1, \ldots, E_m$, then the differential equations\eqref{eq:49} become a system of countably many homogeneous linear equations for these finitely many unknown quantities, with algebraic coefficients. Therefore we may assume that the numerical coefficients of them are integers in $K$, the algebraic number field generated by all coefficients of $E_1, \ldots, E_m$ together.

It is not hard to show that for power series solutions of\eqref{eq:49} the first condition concerning functions of type E, namely, that the coefficients $c_n$ belong to the same algebraic number field $K$ of finite degree, could be weakened to the condition that all coefficients of $E_1, \ldots, E_m$ are algebraic numbers. It could be proved, as a consequence of\eqref{eq:49}, that the field $K$ generated by all coefficients is of finite degree. However, we do not need this property, and we omit the proof.

It is well known that the system\eqref{eq:49} possesses a basis of $m$ solutions $y_{kl}$ ($k=1,\ldots,m$), for $l=1,\ldots,m$. Of course, not all $E_{kl}$ are in general of type E; but all $E_{kl}$ are regular at all finite complex points $x=a$ which are different from the zeros of the polynomial $T(x)$, and the determinant of the $E_{kl}$ does not vanish at $x=a$. Any solution of\eqref{eq:49} takes the form $y_k = c_1 E_{k1} + \dots + c_m E_{km}$ with constant $c_1, \dots, c_m$.

\section{Arithmetical lemmas}

Consider $p$ homogeneous linear equations with $q$ unknown quantities $x_1, \ldots, x_q$ and integral rational coefficients. If $p < q$, they have a non-trivial integral rational solution. It is useful, for different purposes, to obtain an upper estimate of the absolute values of $x_1, \ldots, x_q$ in terms of bounds for the coefficients.

\begin{mylemma}{}{1}
Let
\begin{equation}
y_k = a_{k1}x_1 + \cdots + a_{kq}x_q \quad (k=1,\ldots,p) \tag{50}\label{eq:50}
\end{equation}
be $p$ linear forms with integral rational coefficients and $q$ variables, where $0 < p < q$, and suppose that the absolute values of all $a_{kl}$ are not greater than a given positive rational integer $A$; then there exists a non-trivial integral rational solution $x_1, \ldots, x_q$ of $y_1 = 0, \ldots, y_p = 0$ satisfying the conditions
\[|x_k| < 1 + (qA)^{\frac{p}{q-p}} \quad (k=1,\ldots,q).\]
\end{mylemma}

\begin{proof}
Let $H$ be a positive rational integer and insert in\eqref{eq:50} for $x_1, \ldots, x_q$ independently the $2H+1$ values $0, \pm 1, \ldots, \pm H$. We obtain $(2H+1)^q$ points with integral coordinates $y_1, \ldots, y_p$, all lying in the cube
\[-qAH \le y_k \le qAH \quad (k=1,\ldots,p).\]
Since there are exactly $(2qAH + 1)^p$ different points with integral coordinates in this cube, it follows that at least two different systems $x_1, \dots, x_q$ have the same image point $y_1, \dots, y_p$, if
\begin{equation}
(2qAH+1)^p < (2H+1)^q. \tag{51}\label{eq:51}
\end{equation}
Under this condition we obtain by subtraction a non-trivial integral rational solution $x_1, \dots, x_q$ of $y_1 = 0, \dots, y_p = 0$ satisfying the inequalities
\begin{equation}
|x_k| \le 2H \quad (k=1,\ldots,q). \tag{52}\label{eq:52}
\end{equation}
Now choose for $2H$ the even number in the interval
\[(qA)^{\frac{p}{q-p}} - 1 \le 2H \le (qA)^{\frac{p}{q-p}} + 1\]
of length 2; then\eqref{eq:51} is fulfilled, because of
\[(2qAH+1)^p < (qA)^p (2H+1)^p \le (2H+1)^{q-p}(2H+1)^p = (2H+1)^q\]
and the lemma follows from\eqref{eq:52}.
\end{proof}

Now we generalize\mylem{lem:1} to algebraic number fields $K$.

\begin{mylemma}{}{2}
Suppose that the coefficients of the $p$ linear forms $a_{k1}x_1 + \cdots + a_{kq}x_q$ ($k=1, \ldots, p; p < q$) are integers in $K$, and let $\overline{|a_{kl}|} \le A$; then there exists in $K$ a non-trivial integral solution $x_1, \ldots, x_q$ of $y_1 = 0, \ldots, y_p = 0$ such that
\begin{equation}
\overline{|x_k|} < c + (c q A)^{\frac{p}{q-p}} \tag{53}\label{eq:53}
\end{equation}
where $c$ is a positive constant which only depends upon $K$.
\end{mylemma}

\begin{proof}
Choose a basis $b_1, \ldots, b_h$ of the integers in $K$ relative to the rational number field. Then any integer $\alpha$ in $K$ has the form $\alpha = g_1 b_1 + \cdots + g_h b_h$ with rational integral $g_1, \ldots, g_h$. Solving for $g_1, \dots, g_h$ from the $h$ equations for the conjugates of $\alpha$, it follows that $|g_l| \le \gamma_1 \overline{|\alpha|}$, where $\gamma_1$ only depends upon the choice of the basis. Now write $x_k = x_{k1} b_1 + \cdots + x_{kh} b_h$ with rational integral $x_{k1}, \ldots, x_{kh}$ and express also all $pqh$ products $a_{kl} b_r$ in terms of the basis; then the $p$ equations $y_1 = 0, \dots, y_p = 0$ for $x_1, \ldots, x_q$ become $ph$ homogeneous linear equations for the $qh$ rational integers $x_{11}, \dots, x_{qh}$ with rational integral coefficients of absolute value less than $\gamma_1 \max \overline{|a_{kl} b_r|} \le \gamma_2 A$, where $\gamma_2$ is a positive rational integer depending upon the basis. Applying\mylem{lem:1} we obtain a non-trivial solution satisfying
\[|x_{kl}| < 1 + (\gamma_2 hqA)^{\frac{p}{q-p}} \quad (k=1,\ldots,q; l=1,\ldots,h),\]
and this implies\eqref{eq:53}.
\end{proof}

\section{Approximation forms}

Consider $m$ power series $E_1, \dots, E_m$ and $m$ polynomials $P_1(x), \dots, P_m(x)$ of degree $\le \nu$ with indeterminate coefficients. Obviously, we can choose the $m(\nu+1)$ coefficients, not all 0, such that the approximation form $P_1 E_1 + \dots + P_m E_m$ vanishes at $x=0$ of order $m(\nu+1)-1$ at least. For our further purposes we are interested in the case that $E_1, \dots, E_m$ are of type E; then the coefficients of $P_1, \dots, P_m$ can be taken as integers in the algebraic number field $K$ generated by the coefficients of $E_1, \dots, E_m$. \mylem{lem:2} gives us an upper estimate of the absolute values of the coefficients of $P_1, \dots, P_m$ and their conjugates; however, the estimate in\mylem{lem:2} is useful only in case the ratio $p/q$ is not too near to its upper bound 1. Therefore we now weaken the condition that the approximation form vanishes at $x=0$ of the highest possible order.

\begin{mylemma}{}{3}
Let $E_1, \dots, E_m$ be E-functions with coefficients in $K$, and let an integer $n=1,2,\dots$ be given. There exist $m$ polynomials $P_1(x), \dots, P_m(x)$ of degree $\le 2n-1$ with the following three properties:

1) The coefficients of $P_1, \dots, P_m$ are integers in $K$, not all 0, and the maximum of the absolute values of all their conjugates is $O(n^{(2+\epsilon)n})$, for every given positive $\epsilon$ and $n \to \infty$.

2) The approximation form
\begin{equation}
R = P_1 E_1 + \dots + P_m E_m = \sum_{\nu=0}^{\infty} a_{\nu} \frac{x^{\nu}}{\nu!} \tag{54}\label{eq:54}
\end{equation}
vanishes at $x=0$ of order $(2m-1)n$ at least, so that
\begin{equation}
a_{\nu} = 0 \qquad (\nu = 0, 1, \dots, 2mn-n-1). \tag{55}\label{eq:55}
\end{equation}

3) The coefficients $a_{\nu}$ of $R$ satisfy the condition
\[a_{\nu} = \nu^{\epsilon \nu} O(n^{2n}) \qquad (\nu \ge 2mn-n)\]
uniformly in $\nu$.
\end{mylemma}

\begin{proof}
Put
\[E_{k} = \sum_{\nu=0}^{\infty} c_{k\nu} \frac{x^{\nu}}{\nu!} \qquad (k=1,\dots,m)\]
and
\begin{equation}
P_{k} = (2n-1)! \sum_{\nu=0}^{2n-1} g_{k\nu} \frac{x^{\nu}}{\nu!} \tag{56}\label{eq:56}
\end{equation}
with integral $g_{k\nu}$ in $K$, so that $P_k$ is a polynomial of degree $\le 2n-1$ with integral coefficients in $K$. Computing the coefficients $a_{\nu}$ in the expression\eqref{eq:54} we find
\begin{equation}
P_{k}E_{k} = (2n-1)! \sum_{\nu=0}^{\infty} d_{k\nu} \frac{x^{\nu}}{\nu!}, \tag{57}\label{eq:57}
\end{equation}
\[d_{k\nu} = \sum_{\rho=0}^{2n-1} \binom{\nu}{\rho} g_{k\rho} c_{k,\nu-\rho}\]
\begin{equation}
a_{\nu} = (2n-1)!(d_{1\nu} + \dots + d_{m\nu}) \qquad (\nu = 0, 1, \dots). \tag{58}\label{eq:58}
\end{equation}

The condition\eqref{eq:55} yields $(2m-1)n$ homogeneous linear equations for the $2mn$ unknown quantities $g_{k\rho}$ ($k=1,\dots,m$; $\rho=0,\dots,2n-1$). Determine a positive rational integer $q_n$ such that $q_n c_{k\nu}$ is integral for $k=1,\dots,m$ and $\nu=0,\dots,2mn-n-1$; because of condition 3) in the definition of a function of type E, we may choose $q_n = O(n^{\epsilon n})$. Multiply the equations $a_{\nu}=0$ by $q_n/(2n-1)!$; then the coefficient $q_n \binom{\nu}{\rho} c_{k,\nu-\rho}$ of $g_{k\rho}$ lies in $K$, and all its conjugates again are $O(n^{\epsilon n})$, because of
\begin{equation}
\binom{\nu}{\rho} \le 2^{\nu}, \quad \overline{|c_{k\nu}|} = O(\nu^{\epsilon \nu}) \tag{59}\label{eq:59}
\end{equation}
and $\nu < (2m-1)n$. Applying\mylem{lem:2} with $p=(2m-1)n$, $q=2mn$, $A=O(n^{\epsilon n})$, we find integral $g_{k\nu}$ ($k=1,\dots,m$; $\nu=0,\dots,2n-1$) in $K$, not all 0, such that\eqref{eq:55} is fulfilled and
\[\overline{|g_{k\nu}|} = O(n^{\epsilon n}),\]
because of $\frac{p}{q-p} = 2m-1 = O(1)$. The remaining statements of the lemma readily follow from\eqref{eq:56},\eqref{eq:57},\eqref{eq:58},\eqref{eq:59}, together with the estimate $(2n-1)! = O(n^{2n})$.
\end{proof}

\section{Normal systems}\label{sec:4}

Suppose that the E-functions $E_1, \dots, E_m$ satisfy a system of homogeneous linear differential equations of first order\eqref{eq:49} whose coefficients $Q_{kl}$ are rational functions of $x$, with integral numerical coefficients in the algebraic number field $K$ generated by the coefficients of $E_1, \dots, E_m$. It may happen that the matrix $Q=(Q_{kl})$ decomposes into a number $r$ of quadratic boxes $Q_t=(Q_{kl,t})$, where $t=1,\dots,r$, with $m_1,\dots,m_r$ rows, so that $k,l=1,\dots,m_t$ and $m_1+\dots+m_r=m$. This means that the boxes are arranged in the diagonal of $Q$ and that all elements of $Q$ outside of the boxes are 0. The decomposition is unique if we choose $r$ as large as possible; then we call the $Q_t$ primitive parts of $Q$. Of course, it can happen that $r=1$ and $Q$ itself is primitive.

Corresponding to the decomposition of $Q$ into primitive parts the system\eqref{eq:49} breaks into $r$ separate systems
\begin{equation}
y_{k,t}' = \sum_{l=1}^{m_t} Q_{kl,t}(x)y_{l,t} \qquad (k=1,\dots,m_t; t=1,\dots,r). \tag{60}\label{eq:60}
\end{equation}

Let $y_{l,t} = (y_{kl,t})_{k=1,\dots,m_t}$ for $l=1,\dots,m_t$ be a basis for the solutions of\eqref{eq:60}; then the $m$-rowed matrix $Y$ consisting of the $r$ boxes $Y_t=(y_{kl,t})$ is a solution matrix for\eqref{eq:49}.

Now consider any solution $y_1, \dots, y_m$ of\eqref{eq:49} and introduce the sum
\begin{equation}
R = P_1 y_1 + \dots + P_m y_m \tag{61}\label{eq:61}
\end{equation}
whose coefficients $P_1, \dots, P_m$ are arbitrary polynomials in $x$. We are interested in the case that $R$ vanishes identically in $x$. Using the box decomposition of $Q$ we write $P_{k,t}^*$ ($k=1,\dots,m_t$; $t=1,\dots,r$) instead of $P_1, \dots, P_m$. Expressing the solution $y_1, \dots, y_m$ in terms of the basis we then obtain
\begin{equation}
R = \sum_{k,l,t} P_{k,t}^* c_{l,t} y_{kl,t}, \tag{62}\label{eq:62}
\end{equation}
where $k,l=1,\dots,m_t$ and $t=1,\dots,r$, the $P_{k,t}^*$ being polynomials and the $c_{l,t}$ constant. The sum $R$ trivially equals 0 if all products $P_{k,t}^* c_{l,t}$ are 0, or, in other words, if for each $t=1,\dots,r$ either all polynomials $P_{k,t}^*$ are identically 0 or all constants $c_{l,t}$ are 0. We shall say that the boxes $Y_t$ are independent, if $R$ does not vanish identically in $x$, except in the trivial case.

If the box $Q_t$ is given, then the solution box $Y_t$ only is determined up to a factor $C_t$ to the right, where $C_t$ is an arbitrary $m_t$-rowed non-singular constant matrix; but the passage from $Y_t$ to $Y_t C_t$ does not affect the independence property.

We shall study the algebraical meaning of independence of the solution boxes. Define $R_1 = R$ by\eqref{eq:61} and
\begin{equation}
R_{k+1} = T R_k' \qquad (k=1,2,\dots) \tag{63}\label{eq:63}
\end{equation}
where $T(x)$ is the least common denominator of the $Q_{kl}(x)$. Because of\eqref{eq:49} we can write
\begin{equation}
R_{k} = P_{k1}y_{1} + \dots + P_{km}y_{m} \qquad (k=1,2,\dots), \tag{64}\label{eq:64}
\end{equation}
where
\begin{equation}
P_{k+1,l} = T \left( P_{kl}' + \sum_{g=1}^{m} P_{kg}Q_{gl} \right) \qquad (l=1,\dots,m) \tag{65}\label{eq:65}
\end{equation}
and
\begin{equation}
P_{1l} = P_l. \tag{66}\label{eq:66}
\end{equation}

Clearly, all $P_{kl}$ are polynomials in $x$. Denote by $\Delta = \Delta(x)$ the determinant with the elements $P_{kl}$ ($k,l=1,\dots,m$) and by $\Delta_{kl}$ the minor of $P_{kl}$; then\eqref{eq:64} implies
\begin{equation}
\Delta y_{k} = \sum_{l=1}^{m} \Delta_{kl} R_{l} \qquad (k=1,\dots,m). \tag{67}\label{eq:67}
\end{equation}

There are two cases when we can immediately assure that $\Delta(x)$ vanishes identically in $x$: If in\eqref{eq:62} all $P_{k,t}^*$ for $k=1,\dots,m_t$ and at least one $t$, are identically 0, then the box decomposition of $Q$ shows, because of\eqref{eq:65}, that $m_t$ columns of $\Delta$ are 0; if $R$ is identically 0, but not all $y_1, \dots, y_m$, then\eqref{eq:63} and\eqref{eq:64} imply $\Delta = 0$. Exclude the first case; then the assumption of the second case, because of\eqref{eq:62}, means that the boxes are not independent. Now we prove that, in all other cases, the polynomial $\Delta(x)$ does not vanish identically.

\begin{mylemma}{}{4}
Suppose that the boxes are independent and that, for each $t=1,\dots,r$, not all $P_{k,t}^*$ ($k=1,\dots,m_t$) vanish identically; then $\Delta(x)$ is not identically 0.
\end{mylemma}

\begin{proof}
If $\Delta = 0$, then we could determine $\mu$ polynomials $A_1, \dots, A_{\mu}$ such that
\[A_1 P_{1l} + \dots + A_{\mu} P_{\mu l} = 0 \qquad (l=1,\dots,m), \quad A_{\mu} \neq 0.\]
Let $y_1, \dots, y_m$ be a completely arbitrary solution of\eqref{eq:49} and use the notation\eqref{eq:61} and\eqref{eq:62}. It follows from\eqref{eq:63} and\eqref{eq:64} that
\[A_1 R_1 + \dots + A_{\mu} R_{\mu} = 0\]
\begin{equation}
B_1 R^{(\mu-1)} + B_2 R^{(\mu-2)} + \dots + B_{\mu} R = 0, \tag{68}\label{eq:68}
\end{equation}
where $B_1 = A_{\mu} T^{\mu-1}$ and $B_2, \dots, B_{\mu}$ are polynomials in $x$. Taking for $y_1, \dots, y_m$ the $m$ basic solutions, we see that each of the $m$ functions
\[R_{l,t} = \sum_{k=1}^{m_t} P_{k,t}^* y_{kl,t} \qquad (l=1,\dots,m_t; t=1,\dots,r)\]
satisfies the homogeneous linear differential equation\eqref{eq:68} of order $\mu-1 < m$. Therefore we have a non-trivial homogeneous linear relationship
\begin{equation}
\sum_{l,t} c_{l,t} R_{l,t} = 0 \tag{69}\label{eq:69}
\end{equation}
with constant $c_{l,t}$. But the $r$ boxes are independent, and not all products $P_{k,t}^* c_{l,t}$ ($k,l=1,\dots,m_t$; $t=1,\dots,r$) are 0, because at least one $P_{k,t}^* \neq 0$, for each $t$. So\eqref{eq:69} is impossible.
\end{proof}

We shall say the $m$ functions $E_1, \dots, E_m$ of type E form a normal system if they all are not identically 0 and satisfy the $m$ differential equations\eqref{eq:49}, with rational $Q_{kl}(x)$ having algebraic numerical coefficients, whose solution boxes are independent. This condition of normality will be decisive in the proof of an extension of the result of \S 11, Chapter I, which led to the Lindemann-Weierstrass theorem.

\section{The coefficient matrix of the approximation forms}

From now on we shall assume that the E-functions $E_1, \dots, E_m$ form a normal system and that $P_1, \dots, P_m$ are the polynomials in the approximation form\eqref{eq:54}. Again the polynomials $P_{kl} \ (k = 1, 2, \dots; l = 1, \dots, m)$ will be defined as in\eqref{eq:64} and\eqref{eq:66}, so that the approximation forms
\[R_{k} = P_{k1} E_{1} + \dots + P_{km}E_{m} \quad (k=1,2,\dots)\]
satisfy the equations
\begin{equation}
R_1 = R, \quad R_{k+1} = T R_k' \qquad (70)\tag{70}\label{eq:70}
\end{equation}
None of the functions $E_1, \dots, E_m$ is identically 0. Suppose that the derivatives $E_l^{(k)}(x)$, for $k = 0, \dots, p-1$ and $l = 1, \dots, m$, vanish at $x=0$, but not all $E_l^{(p)}(0) \ (l=1, \dots, m)$. Of course, $p$ may be 0. We denote by $q$ the maximum of the degrees of the $m^2+1$ polynomials $T$ and $T_{kl} \ (k, l = 1, \dots, m)$, and we define
\[t = n + p + q \frac{m(m-1)}{2} - 1.\]

\begin{mylemma}{}{5}
Let $\alpha$ be any complex number different from 0 and the zeros of the polynomial $T(x)$ and suppose that
\begin{equation}
n \ge p + q \frac{m(m-1)}{2} \tag{71}\label{eq:71}
\end{equation}
then the matrix
\begin{equation}
\left( P_{kl}(\alpha) \right)_{\substack{k = 1, \dots, m + t \\ l = 1, \dots, m}} \tag{72}\label{eq:72}
\end{equation}
has the rank $m$.
\end{mylemma}

\begin{proof}
We use the box decomposition of the system\eqref{eq:49} and write again more explicitly $P_{k,t} \ (k=1, \dots, m_t; t= 1, \dots, r)$ instead of $P_1, \dots, P_m$. Not all $P_{k,t}$ are identically 0, because of\mylem{lem:3}. We may choose the notation such that $P_{k,t} = 0$ for $k=1, \dots, m_t$ and $t=\mu+1, \dots, r$; whereas for every $t \le \mu$ at least one $P_{k,t} \neq 0$. Put $m_1 + \dots + m_\mu = \mu$; then $1 \le \mu \le m$. We shall first prove that $\mu = m$.

We apply\mylem{lem:4} with $\mu$ instead of $m$. The assumptions are satisfied, since $E_1, \dots, E_\mu$ form a normal system, a fortiori. Denoting by $\Delta = \Delta(x)$ the determinant of the $P_{kl} \ (k, l = 1, \dots, \mu)$, we see that $\Delta$ does not vanish identically in $x$. By\eqref{eq:67},
\begin{equation}
\Delta y_k = \sum_{l=1}^{\mu} \Delta_{kl} (P_{l1}y_1 + \dots + P_{l\mu} y_{\mu}) \quad (k=1,\dots,\mu) \tag{73}\label{eq:73}
\end{equation}
identically in $x$ and $y_1, \dots, y_\mu$, and in particular
\begin{equation}
\Delta E_{k} = \sum_{l=1}^{\mu} \Delta_{kl} R_{l} \tag{74}\label{eq:74}
\end{equation}
Because of\mylem{lem:3} the power series $R$ vanishes at $x=0$ of order $(2m-1)n$ at least, so that $R_l$, in view of\eqref{eq:70}, vanishes there of order $\ge (2m-1)n - l + 1$. Choose $k$ such that $E_k^{(p)}(0) \neq 0$; then\eqref{eq:74} implies
\begin{equation}
\Delta(x) = x^{(2m-1)n - \mu + 1 - p} \Delta_0(x), \tag{75}\label{eq:75}
\end{equation}
where $\Delta_0(x)$ again is a polynomial and not identically 0. On the other hand, $P_{kl}$ has a degree $\le 2n - 1 + (k-1)q$, by\eqref{eq:65} and\mylem{lem:3}; therefore the $\mu$-rowed determinant $\Delta$ has a degree $\le (2n-1)\mu + q \frac{\mu(\mu-1)}{2}$. If $b$ denotes the degree of $\Delta_0$, then
\begin{equation}
0 \le b \le (2n-1)\mu + q\frac{\mu(\mu-1)}{2} - \left((2m-1)n - \mu + 1 - p\right) \tag{76}\label{eq:76}
\end{equation}
\[= -2n(m-\mu) + n + p + q\frac{\mu(\mu-1)}{2} - 1.\]
Because of the condition\eqref{eq:71}, the right-hand side in\eqref{eq:76} would be negative in case $m - \mu \ge 1$. This proves that $\mu = m$ and
\begin{equation}
b \le t \tag{77}\label{eq:77}
\end{equation}
The number $\alpha$ in the statement of the lemma satisfies $\alpha T(\alpha) \neq 0$. If $\Delta(x)$ vanishes at $x=\alpha$ of order $a$, then, by\eqref{eq:75} and\eqref{eq:77},
\begin{equation}
0 \le a \le t \tag{78}\label{eq:78}
\end{equation}
Consider for a moment the indeterminates $y_1, \dots, y_m$ in\eqref{eq:73} as arbitrary solutions of\eqref{eq:49} and apply $a$ times the operator $TD$. We then obtain the formula
\begin{equation}
\begin{aligned}
& T^{a}(x) \Delta^{(a)}(x)y_{k} + \sum_{l=0}^{a-1} \Delta^{(l)}(x)L_{kl} \\
& \quad = \sum_{l=1}^{m+a} M_{kl}(x) \left(P_{l1}(x)y_{1} + \dots + P_{lm}(x)y_{m}\right) \quad (k=1,\dots,m),
\end{aligned}
\tag{79}\label{eq:79}
\end{equation}
identically in $x$ and the indeterminates $y_1, \dots, y_m$, where the $L_{kl}$ are linear forms in $y_1, \dots, y_m$ whose coefficients are polynomials in $x$, and the $M_{kl}$ are polynomials in $x$. Now insert $x = \alpha$; then
\[T^{a}(\alpha) \Delta^{(a)}(\alpha) = \beta \neq 0\]
and the left-hand side of\eqref{eq:79} reduces to $\beta y_k$. It follows that $y_1, \dots, y_m$ can be expressed as linear combinations of the $m+a$ linear forms $P_{l1}(\alpha)y_1 + \dots + P_{lm}(\alpha)y_m \ (l = 1, \dots, m+a)$. In view of\eqref{eq:78}, this contains the statement of the lemma.
\end{proof}

\section{Estimation of $R_k$ and $P_{kl}$}

The coefficients of $E_1, \dots, E_m$, $T$ and $T Q_{kl}$ ($k, l = 1, \dots, m$) lie in some algebraic number field $K$ of finite degree. By\eqref{eq:65} and\mylem{lem:3}, the coefficients of all polynomials $P_{kl}$ ($k=1, 2, \dots; l=1, \dots, m$) all lie in $K$.

\begin{mylemma}{}{6}
Let $\alpha$ be a number in $K$ and $k \le m+t$, then
\[|R_{k}(\alpha)| = O \left(n^{(3+\epsilon)n - (2m-2)n}\right)\]
\[\overline{|P_{kl}(\alpha)|} = O \left(n^{(3+\epsilon)n}\right) \qquad (l=1,\dots,m).\]
\end{mylemma}

\begin{proof}
If the coefficients of two power series $A=\alpha_0+\alpha_1x+\dots$ and $B=\beta_0+\beta_1x+\dots$ satisfy the conditions $|\alpha_k| \le \beta_k$ for $k=0,1,\dots$, we shall write $A \prec B$. Plainly, $T \prec c(1+x)^q$ and $T Q_{kl} \prec c(1+x)^q$ with some positive constant $c$. We are going to prove by induction that
\begin{equation}
R_{k+1} \prec c^k (1+x)^{kq} \prod_{\nu=0}^{k-1} (\nu q + D) \hat{R} \qquad (k=0,1,\dots) \tag{80}\label{eq:80}
\end{equation}
and
\begin{equation}
P_{k+1,l} \prec c^{k} (1+x)^{kq+2n-1} \prod_{\nu=0}^{k-1} (\nu q + m + 2n-1) O(n^{(2+\epsilon)n}) \qquad (l=1,\dots,m), \tag{81}\label{eq:81}
\end{equation}
where
\begin{equation}
\hat{R} = \sum_{\nu=0}^{\infty} |a_{\nu}| \frac{x^{\nu}}{\nu!} = \sum_{\nu=(2m-1)n}^{\infty} \nu^{\epsilon \nu} \frac{x^{\nu}}{\nu!} O(n^{2n}), \tag{82}\label{eq:82}
\end{equation}
and the estimate\eqref{eq:81} remains valid if the coefficients of $P_{k+1,l}$ are replaced by their conjugates.

Because of\mylem{lem:3}, the statement is true for $k=0$. If it is proved for $k-1 \ge 0$ instead of $k$, we obtain, by\eqref{eq:63} and\eqref{eq:65},
\[
\begin{aligned}
R_{k+1} &= T R'_{k} \prec c(1+x)^{q} c^{k-1} \left\{ (k-1)q(1+x)^{(k-1)q-1} + (1+x)^{(k-1)q}D \right\} \prod_{\nu=0}^{k-2} (\nu q+D) \hat{R} \\
&\prec c^{k} (1+x)^{kq} \prod_{\nu=0}^{k-1} (\nu q+D) \hat{R},
\end{aligned}
\]
\[
\begin{aligned}
P_{k+1,l} &\prec c(1+x)^{q} c^{k-1} \prod_{\nu=0}^{k-2} (\nu q + m + 2n - 1) O(n^{(2+\epsilon)n}) (m+D)(1+x)^{(k-1)q+2n-1} \\
&\prec c^k (1+x)^{kq+2n-1} \prod_{\nu=0}^{k-1} (\nu q + m + 2n-1) O(n^{(2+\epsilon)n}),
\end{aligned}
\]
and this is the required result.

If $k \le m+t = n+O(1)$, then, by\eqref{eq:82},
\[
\begin{aligned}
\prod_{\nu=0}^{k-1} (\nu q + D) \hat{R} &= O(n^{(3+\epsilon)n}) \sum_{\rho=0}^{k} \binom{k}{\rho} \sum_{\nu=(2m-1)n}^{\infty} \nu^{\epsilon\nu} \frac{x^{\nu-\rho}}{(\nu-\rho)!} \\
&\prec O(n^{(3+\epsilon)n}) 2^{k} \sum_{\nu=(2m-1)n-k}^{\infty} (\nu+k)^{\epsilon(\nu+k)} \frac{x^{\nu}}{\nu!} \\
&\prec O(n^{(3+\epsilon)n}) 2^{k} \sum_{\nu=(2m-1)n-k}^{\infty} \{(2\nu)^{2\epsilon\nu} + (2k)^{2\epsilon k}\} \frac{x^{\nu}}{\nu!}
\end{aligned}
\]
\begin{equation}
\left( \prod_{\nu=0}^{k-1} (\nu q + D) \hat{R} \right)_{x=\alpha} = O(n^{(3+\epsilon)n}) O(n^{\epsilon n - (2m-2)n}). \tag{83}\label{eq:83}
\end{equation}
The lemma now follows from\eqref{eq:80},\eqref{eq:81}, and\eqref{eq:83}.
\end{proof}

\section{The rank of $E_1(\alpha), \dots, E_m(\alpha)$}

Let $\omega_1, \dots, \omega_m$ be any complex numbers. We shall say that they have rank $r$ relative to the algebraic number field $K$ if they satisfy $m-r$, and not more than $m-r$, linearly independent homogeneous linear equations
\[\lambda_{k1} \omega_1 + \dots + \lambda_{km} \omega_m = 0 \qquad (k=1,\dots,m-r)\]
with coefficients $\lambda_{kl}$ in $K$. In other words, the set $\omega_1, \dots, \omega_m$ contains $r$ and not more than $r$ numbers which are linearly independent in $K$.

\begin{mylemma}{}{7}
Suppose that $\alpha$ and all coefficients of $E_1, \dots, E_m$ lie in the algebraic number field $K$ of degree $h$. If $E_1, \dots, E_m$ form a normal system and if $\alpha T(\alpha) \neq 0$, then the rank of $E_1(\alpha), \dots, E_m(\alpha)$ relative to $K$ is at least $\frac{m}{2h}$.
\end{mylemma}

\begin{proof}
Since $\alpha$ is a regular point of the coefficients $Q_{kl}$ in\eqref{eq:49}, it follows that not all $E_k(\alpha)=0$ ($k=1,\dots,m$). Suppose that
\begin{equation}
\lambda_{k1} E_1(\alpha) + \dots + \lambda_{km} E_m(\alpha) = 0 \qquad (k=1,\dots,m-r) \tag{84}\label{eq:84}
\end{equation}
are $m-r$ linearly independent equations for $E_1(\alpha), \dots, E_m(\alpha)$ with integral coefficients $\lambda_{kl}$ in $K$. Since not all $E_k(\alpha)=0$, we have $r \ge 1$.

Now apply\mylem{lem:5}, and determine $r$ rows of the matrix\eqref{eq:72}, say for $k=k_1,\dots,k_r$, such that the $r$ linear forms
\begin{equation}
P_{k1}(\alpha) E_1(\alpha) + \dots + P_{km}(\alpha) E_m(\alpha) = R_k(\alpha) \qquad (k=k_1,\dots,k_r) \tag{85}\label{eq:85}
\end{equation}
together with the $m-r$ linear forms in\eqref{eq:84} are independent. Write\eqref{eq:85} before\eqref{eq:84}, denote by $\Lambda$ the determinant of the $m^2$ coefficients $P_{kl}$ and $\lambda_{kl}$, and by $\Lambda_{kl}$ the minor of the element in the $k$-th row and $l$-th column of $\Lambda$; then
\[\Lambda E_{l}(\alpha) = \sum_{g=1}^{r} \Lambda_{gl} R_{k_{g}}(\alpha) \qquad (l=1,\dots,m).\]
Because of\mylem{lem:6},
\[R_{k_{g}}(\alpha) = O(n^{(3+\epsilon)n-(2m-2)n}) \qquad (g=1,\dots,r)\]
\[\overline{|P_{k_g l}(\alpha)|} = O(n^{(3+\epsilon)n}) \qquad (l=1,\dots,m);\]
therefore
\[\Lambda_{gl} = O \left( n^{(3+\epsilon)n(r-1)} \right)\]
\[\Lambda E_{l}(\alpha) = O \left( n^{(3+\epsilon)rn-(2m-2)n} \right)\]
\[\overline{|\Lambda|} = O \left( n^{(3+\epsilon)rn} \right)\]
\[E_{l}(\alpha) \mathfrak{N}(\Lambda) = O \left( n^{(3+\epsilon)rhn-(2m-2)n} \right).\]
Choose a rational positive integer $v$ such that $v \alpha$ is integral; then $v^{2n-1} P_l(\alpha)$ and $v^{2n-1+(k-1)q} P_{kl}(\alpha)$, for $k=1, 2, \dots$ and $l=1, \dots, m$, are integral. Since $\Lambda \neq 0$, it follows that
\[v^{O(n)} \mathfrak{N}(\Lambda) \ge 1.\]
Finally, let $n \to \infty$. The number $E_l(\alpha) \neq 0$ for some $l$; hence
\[3rh \ge 2m-2\]
\[r \ge \frac{2m-2}{3h}.\]
Now the assertion follows, in case $m \ge 3$, because of $h \ge 1$. If $m \ge 3$, we have $\frac{2m-2}{3h} \ge \frac{m}{2h}$, and this means for integers $r, h$ the same as $r \ge \frac{m}{2h}$. In the remaining case $m=1, 2$, the assertion is trivial, because of $r \ge 1$.
\end{proof}

\section{Algebraic independence}

Consider any solution $y_1, \dots, y_m$ of the system\eqref{eq:49}. If $\nu$ is any given rational positive integer, then the number of power products
\[Y = y_1^{\nu_1} y_2^{\nu_2} \dots y_m^{\nu_m},\]
with non-negative rational integral exponents $\nu_1, \dots, \nu_m$ and
\[\nu_1 + \dots + \nu_m \le \nu\]
equals
\[\mu = \mu_{\nu} = \binom{m+\nu}{m}.\]
Because of
\[D \log Y = \sum_{k=1}^{m} \nu_k D \log y_k,\]
the $\mu$ functions $Y$ satisfy a system of $\mu$ homogeneous linear differential equations of first order, whose coefficients are homogeneous linear functions of the $Q_{kl}$ with rational integral numerical coefficients.

\begin{mythm*}{}
Let the E-functions $E_1, \dots, E_m$ be solutions of a system of $m$ homogeneous linear differential equations of first order, whose coefficients are rational functions with algebraic numerical coefficients, and suppose that the $\mu_{\nu}$ power products $E_1^{\nu_1} \dots E_m^{\nu_m}$ ($\nu_1 + \dots + \nu_m \le \nu$) form a normal system, for all $\nu = 1, 2, \dots$. If $\alpha$ is any algebraic number different from 0 and the poles of the $Q_{kl}$, then the $m$ numbers $E_1(\alpha), \dots, E_m(\alpha)$ are not related by an algebraic equation with algebraic coefficients.
\end{mythm*}

\begin{proof}
Let $S(y_1, \dots, y_m)$ be a polynomial in $y_1, \dots, y_m$ whose coefficients are algebraic numbers and not all 0. Denote the total degree of $S$ by $s$, take $\nu \ge s$ and consider the polynomials $y_1^{\nu_1} \dots y_m^{\nu_m} S$ with
\[\nu_1 + \dots + \nu_m \le \nu - s.\]
Their number equals
\[\mu_{\nu-s} = \binom{m-s+\nu}{m}\]
and their total degree is $\le \nu$. If $S$ has the zero $y_1 = E_1(\alpha), \dots, y_m = E_m(\alpha)$, then we obtain $\mu_{\nu-s}$ independent homogeneous linear relations between the $\mu_{\nu}$ numbers $E_1^{\nu_1}(\alpha) \dots E_m^{\nu_m}(\alpha)$ ($\nu_1 + \dots + \nu_m \le \nu$), with coefficients in the algebraic number field $K$ generated by the coefficients of $E_1, \dots, E_m$, $Q_{kl}$, $S$ and the number $\alpha$.

Now apply\mylem{lem:7}, with $\mu_{\nu}$ and the power products $E_1^{\nu_1} \dots E_m^{\nu_m}$ instead of $m$ and $E_1, \dots, E_m$. Because of the normality assumption we infer that
\begin{equation}
\mu_{\nu} - \mu_{\nu-s} \ge \frac{\mu_{\nu}}{2h}, \tag{86}\label{eq:86}
\end{equation}
where $h$ is the degree of the field $K$. But $\mu_{\nu}$ and $\mu_{\nu-s}$ are polynomials of degree $m$ in $\nu$ starting with the same term $\frac{\nu^m}{m!}$; therefore\eqref{eq:86} contains a contradiction as $\nu \to \infty$.

The importance of the theorem lies in the fact that it reduces the arithmetical problem of algebraic independence of the numbers $E_1(\alpha), \dots, E_m(\alpha)$ to the analytical problem of normality of the power products of the functions $E_1(x), \dots, E_m(x)$.

The simplest example is $E_k = e^{\alpha_k x}$ ($k=1, \dots, m$), where $\alpha_1, \dots, \alpha_m$ are algebraic numbers and linearly independent in the rational number field. The $\mu$ power products then take the form $e^{\beta_k x}$ ($k=1, \dots, \mu$) with $\mu$ different algebraic numbers $\beta_k$. The corresponding system\eqref{eq:49} now simply is $y_k' = \alpha_k y_k$ ($k=1, \dots, m$), so that all boxes are one-rowed, and normality means that any equation $P_1 e^{\beta_1 x} + \dots + P_{\mu} e^{\beta_{\mu} x} = 0$ with polynomials $P_k$ implies $P_1 = 0, \dots, P_{\mu} = 0$, which is easily proved. This shows that the Lindemann-Weierstrass theorem is contained in our theorem.
\end{proof}

\section{Hypergeometric E-functions}

Put
\[[\alpha, \nu] = \alpha(\alpha + 1) \dots (\alpha + \nu - 1) \qquad (\nu = 0, 1, \dots)\]
so that $[\alpha, 0] = 1$ and $[\alpha, \nu+1] = (\alpha+\nu)[\alpha, \nu]$. Let $a_1, \dots, a_g$ and $b_1, \dots, b_m$ be rational numbers, $b_k \neq 0, -1, -2, \dots$ ($k=1, \dots, m$), and $m-g=t > 0$. Define
\[c_{n} = \frac{[a_{1},n][a_{2},n]\dots[a_{g},n]}{[b_{1},n][b_{2},n]\dots[b_{m},n]}\]
\[y = \sum_{n=0}^{\infty} c_{n}x^{n} \qquad (n=0,1,\dots)\]
and introduce the operators
\[\Delta_{\lambda} f(x) = D(x^{1+\lambda t} f(x)),\]
\[A = \Delta_{a_1 - a_2} \Delta_{a_2 - a_3} \dots \Delta_{a_{g-1} - a_g} \Delta_{a_g},\]
\[B = \Delta_{b_1 - b_2} \Delta_{b_2 - b_3} \dots \Delta_{b_{m-1} - b_m} \Delta_{b_m - 1},\]
then
\[A x^{tn-1} = (a_1+n)(a_2+n) \dots (a_g+n)t^g x^{t(a_1+n)-1},\]
\[B x^{tn-1} = (b_1+n-1)(b_2+n-1) \dots (b_m+n-1)t^m x^{t(b_1+n-1)-1},\]
so that
\[A \left(\frac{y}{x}\right) = t^g x^{ta_1-1} \sum_{n=0}^{\infty} (a_1+n)(a_2+n) \dots (a_g+n)c_n x^{tn}\]
\[B \left(\frac{y}{x}\right) = t^{m} x^{tb_{1}-1} \sum_{n=-1}^{\infty} (b_{1}+n)(b_{2}+n) \dots (b_{m}+n) c_{n+1} x^{tn}.\]
Since
\[(a_1+n)\dots(a_g+n)c_n = (b_1+n)\dots(b_m+n)c_{n+1} \qquad (n=0,1,\dots),\]
we obtain
\[x^{1-g}(x^{-tb_1}B - t^t x^{-ta_1}A) \left(\frac{y}{x}\right) = t^m(b_1-1)(b_2-1)\dots(b_m-1)x^{-m}.\]
The left-hand side takes the form
\[W = y^{(m)} + Q_1 y^{(m-1)} + \dots + Q_{m-1} y' + Q_m y,\]
where $Q_1, \dots, Q_m$ are polynomials in $x^{-1}$ with rational coefficients. If one of the $m$ numbers $b_1, \dots, b_m$ equals 1, then $y$ satisfies the homogeneous linear differential equation $W=0$ of order $m$. If $b_k \neq 1$ for all $k=1,\dots,m$, we replace $g,m$ by $g+1, m+1$ and define $a_{g+1}=b_{m+1}=1$; this yields a homogeneous linear differential equation $W=0$ of order $m+1$.

Now we shall prove that $y$ is an E-function. Writing
\[y = \sum_{\nu=0}^{\infty} d_{\nu} \frac{x_{\nu}}{\nu!}\]
we have
\begin{equation}
d_{tn} = (tn)! c_{n} = \frac{[a_{1},n]\dots[a_{g},n][1,n]\dots[1,n]}{[b_{1},n]\dots[b_{g},n][b_{g+1},n]\dots[b_{m},n]} \cdot \frac{(tn)!}{(n!)^{t}} \qquad (n=0,1,\dots) \tag{87}\label{eq:87}
\end{equation}
and $d_{\nu} = 0$ otherwise. It is clear that
\[\frac{[a,n]}{[b,n]} = \prod_{k=1}^{n} \left\{ \left(1 + \frac{a-1}{k}\right) \Big/ \left(1 + \frac{b-1}{k}\right) \right\} = O(e^{\epsilon n}) \qquad (b \neq 0, -1, -2, \dots),\]
for any given $\epsilon > 0$ and $n \to \infty$, and
\[\frac{(tn)!}{(n!)^{t}} \le (1+\dots+1)^{tn} = t^{nt};\]
hence
\[d_n = O(n^{\epsilon n}).\]

It remains to prove that the third condition in the definition of an E-function is fulfilled, namely, that the least common denominator of the rational numbers $d_0,\dots,d_n$ is $O(n^{\epsilon n})$. Since the number $(tn)!/(n!)^t$ is integral, it suffices to prove, because of\eqref{eq:87}, that the least common denominator of the $n+1$ numbers $[a,k] / [b,k]$ ($k=0,\dots,n$) is $O(n^{\epsilon n})$, for any rational $a,b$ and $b \neq 0,-1,-2,\dots$. Put $a=\alpha/\beta$, $b=\gamma/\delta$, where $(\alpha,\beta)=1$, $(\gamma,\delta)=1, \delta>0$, then
\[\frac{\beta^{2k}[a,k]}{\delta^{k}[b,k]} = \frac{\beta^{k}\alpha(\alpha+\beta) (\alpha+2\beta)\dots(\alpha+(k-1)\beta)}{\gamma(\gamma+\delta) (\gamma+2\delta)\dots(\gamma+(k-1)\delta)} = \frac{M_{k}}{N_{k}} \qquad (k=0,\dots,n),\]
say. Consider any prime factor $p$ of $N_k$; then $(p,\delta)=1$. If $\nu$ runs over $p^l$ consecutive rational integers, then one and only one of the $p^l$ corresponding integers $\gamma+\nu\delta$ is divisible by $p^l$, for $l=1,2,\dots$. Therefore at least $[kp^{-l}]$ and at most $1+[kp^{-l}]$ of the $k$ factors $\gamma,\gamma+\delta,\dots,\gamma+(k-1)\delta$ in $N_k$ are divisible by $p^l$. In case $p^l > |\gamma|+(k-1)\delta$ none of these factors is divisible by $p^l$. Suppose that $N_k$ is divisible by $p^s$, but not by $p^{s+1}$, then
\[\sum_{l} [kp^{-l}] \le s \le \sum_{l} (1+[kp^{-l}]),\]
where $l=1,2,\dots$ is restricted by the condition
\[1 \le \frac{\log (|\gamma| + (k-1)\delta)}{\log p}.\]
Hence
\begin{equation}
\left[\frac{k}{p}\right] \le s \le \left[\frac{k}{p}\right] + O\left(\frac{k}{p^{2}}\right) + O\left(\frac{\log(k+1)}{\log p}\right) \tag{88}\label{eq:88}
\end{equation}
and
\[p \le |\gamma| + (k-1)\delta.\]
In case $(p,\beta)=1$ the left-hand estimate in\eqref{eq:88} shows that also the numerator is divisible by $p^{[k/p]}$, and this remains true for the prime factors of $\beta$, because of the factor $\beta^k$ in $M_k$. Denote by $r_p$ the exponent of $p$ in the reduced denominator of the number $M_k/N_k$; then
\[r_{p} = O\left(\frac{n}{p^{2}}\right) + O\left(\frac{\log n}{\log p}\right),\]
for $k=0,1,\dots,n$. Therefore the least common denominator $q_n$ of $[a,k] / [b,k]$ ($k=0,\dots,n$) satisfies
\[\log q_n = O(n) + \sum_{p=O(n)} \left\{ O\left(\frac{n}{p^2}\right) + O\left(\frac{\log n}{\log p}\right) \right\} \log p = O(n) + \pi(O(n)) O(\log n) = O(n).\]
Here we used the elementary upper estimate $\pi(x) = O\left(\frac{x}{\log x}\right)$ for the prime number function. This completes the proof.

Since $\sum_{n=0}^{\infty} c_n x^n$ for $c_n = \frac{[\alpha, n][\beta, n]}{[\gamma, n][1, n]}$ is the hypergeometric series, we shall speak of the functions $y$ defined in this section as hypergeometric E-functions. Performing the substitution $x \to \lambda x$ for arbitrary algebraic $\lambda$ and taking any polynomial in $x$ and finitely many hypergeometric E-functions, with algebraic coefficients, we get again an E-function satisfying a homogeneous linear differential equation whose coefficients are rational functions of $x$. It would be interesting to find out whether all such E-functions can be constructed in the preceding manner.

\section{The Bessel differential equation}\label{sec:10}

We shall apply our theorem to the particular function
\begin{equation}
\mathbf{K} = \mathbf{K}_{\lambda} (\mathbf{x}) = \sum_{n=0}^{\infty} \frac{(-1)^n}{n! (\lambda + 1)(\lambda + 2) \dots (\lambda + n)} \left(\frac{\mathbf{x}}{2}\right)^{2n} \quad (\lambda \neq -1, -2, \dots) \tag{89}\label{eq:89}
\end{equation}
This is a special case of the hypergeometric E-functions introduced in \S 9: Choose $q=0, m=2, b_1=1, b_2=\lambda+1$ and substitute $-\frac{x^2}{4}$ for $x$. The differential equation for $K$ then becomes
\[K'' + \frac{2\lambda + 1}{x} K' + K = 0,\]
so that the two E-functions $y_1 = K, y_2 = K'$ are solutions of the system
\[y_1' = y_2, \quad y_2' = -y_1 - \frac{2\lambda + 1}{x} y_2.\]
We have to investigate whether these functions satisfy the normality condition in the theorem. This requires a lengthy discussion which, however, is not without interest in itself. The answer will be given in \S 13, and it will turn out that the normality condition is satisfied for all rational $\lambda \neq \pm \frac{1}{2}, \pm \frac{3}{2}, \dots.$

It is practical to consider the differential equation for $y = x^\lambda K$ instead of $K$. This gives the Bessel differential equation
\begin{equation}
y'' + \frac{1}{x} y' + \left(1 - \frac{\lambda^2}{x^2}\right)y = 0 \tag{90}\label{eq:90}
\end{equation}
with the particular solution
\begin{equation}
J_{\lambda}(x) = \frac{1}{\Gamma(\lambda+1)} \left(\frac{x}{2}\right)^{\lambda} K_{\lambda}(x), \tag{91}\label{eq:91}
\end{equation}
the Bessel function of index $\lambda$. For our next purposes, the parameter $\lambda$ in\eqref{eq:90} may be any complex number, not necessarily rational. Let any solution $y$ of\eqref{eq:90} be given, not identically 0. This function is everywhere regular, with possible exception of the points $0$ and $\infty$. It is our first aim to prove that the three functions $x, y, y'$ are not related by an algebraic equation with constant coefficients, except when $2\lambda$ is odd.

We start by proving the rather trivial statement that $y$ is not an algebraic function of $x$. Otherwise there would exist an expansion
\[y = \sum_{k=0}^{\infty} c_k x^{r_k}\]
with decreasing rational exponents $r_0 > r_1 > \dots$ and $c_0 \neq 0$, converging near $x=\infty$. Inserting into\eqref{eq:90} we obtain the contradiction $c_0=0$.

Now consider more generally an arbitrary homogeneous linear differential equation of the second order
\begin{equation}
w'' + A(x)w' + B(x)w = 0 \tag{92}\label{eq:92}
\end{equation}
whose coefficients $A$ and $B$ belong to a given field $L$ of analytic functions of $x$. We shall assume that $L$ is closed with respect to differentiation, in other words, that $L$ contains the derivatives of all its elements. For instance, $L$ may be the field of rational functions of $x$. Now suppose that\eqref{eq:92} has a particular solution $w_0$, which is not algebraic over $L$, but satisfies an algebraic differential equation of the first order, with coefficients in $L$. We shall prove that then there exists also a solution of\eqref{eq:92} whose logarithmic derivative is algebraic over $L$. Our assumption means that one can find a polynomial $P(y,z)$ in two indeterminates $y,z$, with $P(w_0, w_0') = 0$ identically in $x$; moreover, $P$ is not independent of $z$ itself.

If $Q(y,z)$ is any polynomial in $y,z$ with coefficients in $L$, we define
\begin{equation}
Q^*(y,z) = Q_x + zQ_y - (Az+By)Q_z; \tag{93}\label{eq:93}
\end{equation}
then $Q^*(y,z)$ again is a polynomial in $y,z$ with coefficients in $L$ and
\begin{equation}
\frac{d}{dx} Q(w,w') = Q^*(w,w') \tag{94}\label{eq:94}
\end{equation}
for every solution $w$ of\eqref{eq:92}. We may assume that $P(y,z)$ is irreducible. Consider $P(y,z)$ and $P^*(y,z)$ as polynomials in $z$ alone and introduce their resultant $R(y)$; this is a polynomial in $y$ with coefficients in $L$. Because of\eqref{eq:94}, the differential equation $P(w_0,w_0')=0$ implies $P^*(w_0,w_0')=0$, hence $R(w_0)=0$. But $w_0$ is not algebraic over $L$, so that $R(y)$ vanishes identically in $y$. This proves that $P$ and $P^*$, as polynomials in $z$, are not coprime. Since $P$ is irreducible, we obtain
\begin{equation}
P^*(y,z) = T(y,z) P(y,z) \tag{95}\label{eq:95}
\end{equation}
where $T$ is a polynomial in $y,z$ with coefficients in $L$. Now take in $P(y,z)$ the aggregate $H(y,z)$ of terms of highest total degree in $y,z$, so that $H$ is a homogeneous polynomial in $y,z$ of positive degree $t$, with coefficients in $L$ and not all 0. The definition\eqref{eq:93} shows that $H^*$ again is a homogeneous polynomial in $y,z$ of degree $t$, and that $P^*-H^*=(P-H)^*$ is of total degree $< t$. Comparing degrees in\eqref{eq:95} we see that $T$ cannot have positive total degree; therefore $T$ is independent of $y,z$ and a function in $L$; moreover
\begin{equation}
H^*(y,z) = T H(y,z). \tag{96}\label{eq:96}
\end{equation}

Consider the homogeneous linear differential equation
\begin{equation}
v' = Tv \tag{97}\label{eq:97}
\end{equation}
of the first order. Its general solution is $v= c v_0$ where $v_0 \neq 0$ is a particular solution and $c$ an arbitrary constant. Because of\eqref{eq:94} and\eqref{eq:96}, the function $H(w,w')$ is a solution of\eqref{eq:97} for every w satisfying\eqref{eq:92}. Choose two independent solutions $w_1, w_2$ of\eqref{eq:92}, then $w= \lambda_1 w_1 + \lambda_2 w_2$ is the general solution of\eqref{eq:92}, with arbitrary constants $\lambda_1, \lambda_2$. Therefore
\[H(\lambda_1 w_1 + \lambda_2 w_2, \lambda_1 w_1' + \lambda_2 w_2') = c(\lambda_1, \lambda_2) v_0\]
where $c(\lambda_1, \lambda_2)$ is independent of $x$. But the left-hand side is a homogeneous polynomial of degree $t$ in $\lambda_1, \lambda_2$; hence
\[c(\lambda_1, \lambda_2) = c_0 \lambda_1^{t} + c_1 \lambda_1^{t-1} \lambda_2 + \dots + c_t \lambda_2^{t}\]
with constant $c_0, \dots, c_t$.

Finally determine two constant values $\lambda_1, \lambda_2$, not both 0, such that $c(\lambda_1, \lambda_2)=0$. Then the corresponding particular solution $w=\lambda_1 w_1 + \lambda_2 w_2$ of\eqref{eq:92} satisfies the differential equation
\[H\left(1,\frac{w'}{w}\right) = 0.\]
This means that the logarithmic derivative of $w$ is algebraic over $L$.

\section{Determination of the exceptional case}\label{sec:11}

We apply the result of \autoref{sec:10} to the Bessel differential equation\eqref{eq:90} and take for $L$ the field of rational functions of $x$. Suppose that $y_0 \neq 0$ is a solution of\eqref{eq:90} which satisfies an algebraic differential equation $P(y_0,y'_0)=0$ whose coefficients are polynomials in $x$ and not all 0. It follows from our result that then\eqref{eq:90} has another particular solution $y \neq 0$ whose logarithmic derivative $\frac{y'}{y} = u$ is an algebraic function of $x$. The function $y$ is regular for all $x \neq 0, \infty$; therefore the only possible branch points of $u$ lie at $0$ and $\infty$. We shall prove that no branch of $u$ is ramified at $\infty$. This implies that also $0$ is no branch point, and $u$ has to be a rational function of $x$.

Let
\begin{equation}
u = \sum_{k=0}^{\infty} c_k x^{r_k} \tag{98}\label{eq:98}
\end{equation}
be the power series expansion of any branch of $u$ near $x=\infty$, with rational exponents $r_0 > r_1 > \dots$ and $c_0 \neq 0$. Because of\eqref{eq:90}, the function $u$ satisfies the special Riccati differential equation
\[u' + u^2 + \frac{1}{x}u = \frac{\lambda^2}{x^2} - 1;\]
hence
\begin{equation}
\sum_{k=0}^{\infty} (r_k+1)c_k x^{r_k-1} + \sum_{k,l=0}^{\infty} c_k c_l x^{r_k+r_l} = \frac{\lambda^2}{x^2} - 1. \tag{99}\label{eq:99}
\end{equation}
Comparing coefficients we obtain
\[r_0 = 0, \quad c_0 = \pm i.\]
For any $n=1,2,\dots$ the two terms corresponding to $k=0, l=n$ and $k=n, l=0$ in the double sum give the contribution $2 c_0 c_n x^{r_n}$ to the left-hand side of\eqref{eq:99}. If $c_n \neq 0$, the exponent $r_n$ has to equal one of the exponents $r_k-1 \ (k < n)$, $r_k+r_l \ (k, l < n)$ and $-2$. We infer by induction that we may restrict the exponents to the sequence of rational integers $-k \ (k=0,1,\dots)$. In particular, $r_1 = -1$ and
\[2c_0c_1 + c_0 = 0, \quad c_1 = -\frac{1}{2}.\]
Therefore
\begin{equation}
u = \pm i - \frac{1}{2x} + \dots \tag{100}\label{eq:100}
\end{equation}
is regular and unramified at $\infty$.

Now we know that $u$ is a rational function of $x$. Apply\eqref{eq:98} and\eqref{eq:99} to the power series expansion of $u$ near $x=0$, so that the exponents $r_k$ are consecutive integers in ascending order. It follows that this power series takes the form
\begin{equation}
u = \pm \frac{\lambda}{x} + \dots \tag{101}\label{eq:101}
\end{equation}
If $x_0 \neq 0, \infty$ is a zero of $y$, of order $a \ge 1$, then $u$ has at $x=x_0$ a pole of first order with the residue $a$. Since $u$ is a rational function, $y$ has only finitely many zeros $\neq 0, \infty$, say $x_1, \dots, x_h$, taken with their multiplicity. The function $u$ is regular at all other points $\neq 0$. It follows from\eqref{eq:100} and\eqref{eq:101} that
\[u = \pm i \pm \frac{\lambda}{x} + \sum_{k=1}^{h} \frac{1}{x-x_k}\]
is the partial fraction expansion of $u$ and that
\[h \pm \lambda = -\frac{1}{2}.\]
Therefore $2\lambda = \pm (2h+1)$ is an odd number.

We have reached the first aim announced in \autoref{sec:10}, namely, we have proved that no solution $y \neq 0$ of the Bessel differential equation satisfies an algebraic differential equation of first order whose coefficients are polynomials in $x$, provided $2\lambda$ is not an odd number. It is known that the case $2\lambda = \pm (2h+1)$ really is exceptional. For any given $h=0,1,\dots$, the functions
\[y_{1} = x^{h+\frac{1}{2}} \frac{d^{h}}{d(x^{2})^{h}} \frac{e^{ix}}{x},\]
\[y_{2} = x^{h+\frac{1}{2}} \frac{d^{h}}{d(x^{2})^{h}} \frac{e^{-ix}}{x}\]
are two independent solutions of\eqref{eq:90}, with $\lambda = \pm (h + \frac{1}{2})$. Every solution then takes the form
\[y = x^{-h-\frac{1}{2}} (Ae^{ix} + Be^{-ix})\]
where $A(x)$ and $B(x)$ are certain polynomials. Computing $y^2$, $yy'$, $(y')^2$ and eliminating $e^{2ix}$, $e^{-2ix}$ one obtains the quadratic differential equation of first order
\[C_1 y^2 + C_2 yy' + C_3 (y')^2 = C_4\]
where $C_1, C_2, C_3, C_4$ are polynomials in $x$ and not all 0. From now on we shall exclude the exceptional case.

\section{Algebraic relations involving different Bessel functions}\label{sec:12}

Let two independent solutions $y_1$, $y_2$ of the Bessel differential equation be given. Then the function $y_1 y_2' - y_2 y_1' = \Delta$ satisfies $\Delta' = -\frac{\Delta}{x}$, whence
\begin{equation}
y_1 y_2' - y_2 y_1' = \frac{c}{x} \tag{102}\label{eq:102}
\end{equation}
with constant $c \neq 0$. This shows that the five functions $y_1, y_1', y_2, y_2'$ and $x$ are algebraically dependent.

We shall prove that the four functions $y_1, y_1', y_2, x$ are algebraically independent. Because of the result of \autoref{sec:11} we have to prove that $y_2$ is not algebraic over the field $M$ of the rational functions of $y_1, y_1', x$. Suppose that there exists an irreducible polynomial
\[P(t) = t^n + \dots\]
with coefficients in $M$, such that
\[P(y_2) = 0\]
identically in $x$. We define
\begin{equation}
P^*(t) = \left(\frac{y_1'}{y_1} t + \frac{c}{x y_1}\right) P_t(t) + \frac{dP(t)}{dx} \tag{103}\label{eq:103}
\end{equation}
\[= n \frac{y_1'}{y_1} t^n + \dots;\]
this again is a polynomial of degree $n$ in $t$, with coefficients in $M$. For arbitrary constant $\lambda$, the function $y_0 = y_2 + \lambda y_1$ is a solution of the linear differential equation of first order
\[y_0' = \frac{y_1'}{y_1} y_0 + \frac{c}{x y_1}\]
for $y_0$. The definition\eqref{eq:103} implies
\begin{equation}
P^*(y_0) = \frac{dP(y_0)}{dx} \tag{104}\label{eq:104}
\end{equation}
and, in particular,
\[P^*(y_2) = 0.\]
It follows that $P^*(t)$ is divisible by $P(t)$, whence
\[P^*(t) = n \frac{y_1'}{y_1} P(t)\]
and, because of\eqref{eq:104},
\begin{equation}
P(y_0) = b y_1^n \tag{105}\label{eq:105}
\end{equation}
where $b$ is independent of $x$. The function $P(y_2 + \lambda y_1)$ is a polynomial in $\lambda$ of degree $n$; therefore $b = b_0 + b_1 \lambda + \dots + b_n \lambda^n$ with constant $b_0, \dots, b_n$, and
\[P_t(y_2) = b_1 y_1^{n-1}.\]
This is an algebraic equation for $y_2$ of degree $n-1$, with coefficients in $M$; hence $n=1$ and $y_2$ itself lies in $M$, so that
\[y_2 = \frac{f}{g}\]
where $f$ and $g$ are polynomials in $y_1, y_1'$ and $x$. Take in $f$ and $g$ the homogeneous aggregates of highest total degree in $y_1, y_1'$, say $f_0$ of degree $\phi$ and $g_0$ of degree $\gamma$, and introduce the difference $\delta = \phi - \gamma$ as total degree of the ratio $f/g$.

Put $f_0/g_0 = v$; then $v$ is homogeneous of degree $\delta$ in $y_1, y_1'$, and the difference $f/g - v$ has total degree $< \delta$. Because of the differential equation\eqref{eq:90}, the same is true for $v'$ and $(f/g)' - v'$. Inserting\eqref{eq:105} into\eqref{eq:102} we obtain an identity in $y_1, y_1', x$. On the left-hand side of\eqref{eq:102} the terms of highest total degree yield $y_1 v' - v y_1'$ of degree $\delta + 1$, whereas the right-hand side has total degree 0. Therefore $\delta \ge -1$.

Suppose first that $\delta > -1$; then $y_1 v' - v y_1' = 0$, whence $v = c_1 y_1$ with constant $c_1$, and $\delta = 1$. Substituting $y_2 + c_1 y_1$ for $y_2$, we have to replace $f/g$ by $f/g + v$, and the new $f/g$ will have total degree $< 1$. This leaves us with the case $\delta = -1$ and
\[y_1 v' - v y_1' = \frac{c}{x},\]
so that $v$ is a rational function of $x, y_1, y_1'$ and homogeneous of degree $-1$ in $y_1, y_1'$ which satisfies the Bessel differential equation. Write $v = v(y_1, y_1')$. Using once more the algebraic independence of $x, y_1, y_1'$ we see that also $v(\lambda_1 y_1 + \lambda_2 y_2, \lambda_1 y_1' + \lambda_2 y_2')$, for arbitrary constant $\lambda_1$ and $\lambda_2$, is a solution of\eqref{eq:90}. Hence
\begin{equation}
v(\lambda_1 y_1 + \lambda_2 y_2, \lambda_1 y_1' + \lambda_2 y_2') = \Lambda_1 y_1 + \Lambda_2 y_2, \tag{106}\label{eq:106}
\end{equation}
where $\Lambda_1, \Lambda_2$ are independent of $x$. Since
\[v y_2' - y_2 v' = \frac{c \Lambda_1}{x}, \quad y_1 v' - v y_1' = \frac{c \Lambda_2}{x},\]
it follows that $\Lambda_1$ and $\Lambda_2$ are rational functions of $\lambda_1$ and $\lambda_2$, homogeneous of dimension $-1$, not both identically 0, and with constant coefficients. Let $\Lambda_0$ be the least common denominator of $\Lambda_1$ and $\Lambda_2$; this is a homogeneous polynomial of degree $\ge 1$. Now multiply\eqref{eq:106} by $\Lambda_0 g_0$ and choose for $\lambda_1, \lambda_2$ two numbers, not both 0, such that $\Lambda_0 = 0$. It follows that the particular solution $y = \lambda_1 y_1 + \lambda_2 y_2$ of\eqref{eq:90} satisfies the algebraic differential equation $g_0(y, y') = 0$, which is impossible.

\section{Normality --- Bessel functions}\label{sec:13}

Now we are going to prove finally that the normality condition of the theorem is fulfilled in case of the two E-functions $E_1 = K_\lambda(x), E_2 = K_\lambda'(x)$, where $K_\lambda(x)$ is defined in\eqref{eq:89} and $\lambda$ denotes a rational number $\neq -1, -2, \dots$. We have $E_1' = E_2, E_2' = -E_1 - \frac{2\lambda+1}{x} E_2$; therefore the $q+1$ functions $z_{kq} = \frac{1}{k!} E_1^{q-k} E_2^k \ (k=0,\dots,q)$ satisfy for every given $q=0,1,2,\dots$ the system of differential equations
\begin{equation}
z_{kq}' = k z_{k-1, q} - (q-k) \frac{2\lambda+1}{x} z_{kq} - (q-k) z_{k+1, q} \tag{107}\label{eq:107}
\end{equation}
\[(k=0,\dots,q)\]
where $z_{-1, q}$ and $z_{q+1, q}$ mean 0. Our problem is to prove that the $\frac{(\nu+1)(\nu+2)}{2}$ functions $z_{kq} \ (k=0,\dots,q; q=0,\dots,\nu)$ constitute for every $\nu=0,1,2,\dots$ a normal system.

If $z$ is any solution of
\begin{equation}
Z'' + \frac{2\lambda+1}{x} Z' + Z = 0 \tag{108}\label{eq:108}
\end{equation}
then the $q+1$ functions $\frac{1}{k!} z^{q-k} (z')^k \ (k=0,\dots,q)$ form a solution of\eqref{eq:107}. Take $z = \rho_1 z_1 + \rho_2 z_2$, where $z_1, z_2$ are two independent solutions of\eqref{eq:108} and $\rho_1, \rho_2$ arbitrary constants, and define
\begin{equation}
\frac{1}{k!} z^{q-k} (z')^k = \sum_{l=0}^q \rho_1^l \rho_2^{q-l} \psi_{kl,q}; \tag{109}\label{eq:109}
\end{equation}
then the $q+1$ functions $\psi_{kl,q} \ (k=0,\dots,q)$ constitute for each $l=0,\dots,q$ a solution of\eqref{eq:107}. Since $\psi_{0l,q} = \binom{q}{l} z_1^l z_2^{q-l}$, these $q+1$ solutions are linearly independent, because otherwise one would obtain a homogeneous algebraic equation of degree $q$ for $z_1, z_2$, with constant coefficients not all 0, in contradiction to the independence of $z_1$ and $z_2$.

This shows that $\psi_{kl,q}$ is a polynomial in $y_1, y_1^{-1}, y_1', y_2'$ whose coefficients are algebraic functions of $x$, since $\lambda$ is a rational number.

Now suppose that $R=0$, identically in $x$. The result of \autoref{sec:12} implies that then $R$ vanishes identically in $x, y_1, y_1', y_2$. If all products $P_{kq} c_{lq} = 0$ for $k,l=0,\dots,q$ and $q=\mu+1,\dots,\nu$, we consider in $R$ the terms of total degree $\mu$ in $y_1, y_1', y_2$. Because of\eqref{eq:109}, we obtain from $\psi_{kl,\mu}$ the contribution $\binom{\mu}{l} x^{-\lambda q} \left(\frac{y_1'}{y_1} - \frac{\lambda}{x}\right)^{\mu-k} y_1^l y_2^{\mu-l}$. Therefore
\[\sum_{0 \le k, l \le \mu} P_{k\mu} c_{l\mu} \binom{\mu}{l} \left( \frac{y_1'}{y_1} - \frac{\lambda}{x} \right)^{\mu - k} \left( \frac{y_2}{y_1} \right)^{\mu - l} = 0\]
identically in $x, y_1, y_1', y_2$, and this implies $P_{k\mu} c_{l\mu} = 0 \ (k,l=0,\dots,\mu)$. It follows that all $P_{kq} c_{lq}$ are 0, and the proof is complete.

The only singularity of the coefficients in the differential equation\eqref{eq:108}, for finite $x$, is $x=0$. Let $\alpha$ be any algebraic number $\neq 0$, and $\lambda$ a rational number $\neq \pm \frac{1}{2}, -1, \pm \frac{3}{2}, -2, \dots$. It follows from the theorem that the values $K_\lambda(\alpha)$ and $K_\lambda'(\alpha)$ are not related by any algebraic equation with algebraic coefficients. In particular, $K_\lambda(\alpha)$ itself is transcendental. This implies that all zeros of $K_\lambda(\alpha)$ are transcendental, and, because of\eqref{eq:91}, the same holds for the zeros $\neq 0$ of the Bessel function $J_\lambda(x)$.

Putting
\[w = f_\lambda(x) = \sum_{n=0}^\infty \frac{x^n}{n!(\lambda+1)\dots(\lambda+n)}\]
we have
\[K_\lambda(x) = f_\lambda\left(-\frac{x^2}{4}\right)\]
and the differential equation
\[x w'' + (\lambda + 1) w' = w,\]
which leads to the continued fraction
\[\frac{w}{w'} = \lambda + 1 + \frac{x}{\lambda + 2 + \frac{x}{\lambda + 3 + \dots}}\]
Since
\[\frac{w'}{w} = \frac{-1}{\sqrt{-x}} \frac{K_\lambda'(2\sqrt{-x})}{K_\lambda(2\sqrt{-x})},\]
it follows that the value of this continued fraction is transcendental for all algebraic $x \neq 0$ and all rational $\lambda$. In this statement the values $\lambda = \pm \frac{1}{2}, \pm \frac{3}{2}, \dots$ are included, because the remark at the end of \autoref{sec:11} shows that in this case the transcendency follows from that of $e^x$.

In particular, the number
\[1 + \frac{1}{2 + \frac{1}{3 + \dots}}\]
is transcendental.

\section{Additional remarks}

The preceding example of the Bessel differential equation makes clear that the application of the theorem to a given system of linear differential equations of first order requires a detailed study of algebraical and analytical properties of the solutions. Of course, it may happen that the normality condition is not satisfied. The simplest example is provided by the special case $\lambda = -\frac{1}{2}$ in\eqref{eq:89}; then $K = \cos x, K' = -\sin x, K'' = -K$, and the solution box
\[Y = \begin{pmatrix} \cos x & \sin x \\ -\sin x & \cos x \end{pmatrix}\]
is not independent in the sense of \S 4, since $\cos^2 x + \sin^2 x - 1 = 0$. The inner reason is that the substitution $z_1 = y_1 - i y_2, z_2 = y_1 + i y_2$ transforms the system $y_1' = y_2, y_2' = -y_1$ into $z_1' = i z_1, z_2' = -i z_2$, and the transformed $Y$ breaks up into the two boxes $e^{ix}, e^{-ix}$. This makes it probable that the independence of the solution boxes in \S 4 can be expressed as a property of the ring $P$ of linear transformations
\[y_k = A_{k1} y_1 + \dots + A_{km} y_m \quad (k=1,\dots,m),\]
whose coefficients are rational functions of $x$, and which map the linear set of solutions of the differential equations\eqref{eq:49} into itself. It is not difficult to find this connection if one uses the results of A. Loewy on homogeneous linear differential equations. However, we did not need $P$ for the application to Bessel functions, so we only mention it here.

Our result concerning $K_\lambda(\alpha)$ can be generalized in the following way. We consider the $m=2r$ functions $K_\lambda(\alpha_l x), K_\lambda'(\alpha_l x) \ (l=1,\dots,r)$, where $\alpha_1, \dots, \alpha_r$ are given algebraic numbers whose squares all are different from each other and from $0$. It is possible to prove that the condition of the theorem is satisfied, if $2\lambda$ is not odd; therefore the numbers $K_\lambda(\alpha_1), K_\lambda'(\alpha_1), \dots, K_\lambda(\alpha_r), K_\lambda'(\alpha_r)$ are not related by an algebraic equation with algebraic coefficients. This again is a special case of the following statement: Let $\lambda_1, \dots, \lambda_s$ be rational numbers, different from $-1, \pm \frac{1}{2}, -2, \pm \frac{3}{2}, \dots$, and let none of the numbers $\pm \lambda_k \pm \lambda_l \ (1 \le k < l \le s)$ be integral; then the $2rs$ numbers $K_\lambda(\alpha), K_\lambda'(\alpha) \ (\lambda = \lambda_1, \dots, \lambda_s; \alpha = \alpha_1, \dots, \alpha_r)$ are algebraically independent. The proof has not been carried through in detail and we leave it as an interesting exercise.

Another exercise is provided by the hypergeometric E-function
\[y = 1 + \frac{\kappa}{\lambda} \cdot \frac{x}{1!} + \frac{\kappa(\kappa+1)}{\lambda(\lambda+1)} \cdot \frac{x^2}{2!} + \dots = \int_0^1 t^{\kappa-1} (1-t)^{\lambda-\kappa-1} e^{tx} dt \Big/ \int_0^1 t^{\kappa-1} (1-t)^{\lambda-\kappa-1} dt\]
with rational $\kappa, \lambda \neq 0, -1, -2, \dots$, which is a solution of
\[x y'' + (\lambda - x) y' = \kappa y.\]
It seems more difficult to carry over the investigations of \autoref{sec:10}, \autoref{sec:11}, \autoref{sec:12}, \autoref{sec:13} to linear differential equations of order greater than two, but it would be worthwhile to try.
